\chapter{Aide-mémoire des commandes}

\section*{git}
\begin{terminal}
# ce qui a changé depuis le dernier commit
git status
# enregistrer une idée
git add fichier && git commit -m "…"
# les vingt derniers commits
git log --oneline -20
# quand ce mot est apparu ou a disparu
git log -S "mot" -- fichier
# le commit de chaque ligne
git blame fichier
# figer une version
git tag nom && git push origin nom
# ce qui a changé depuis
git diff nom -- fichier
\end{terminal}

\section*{Tickets (GitHub, en ligne de commande)}
\begin{terminal}
# les tickets ouverts
gh issue list
# lire un ticket
gh issue view 95
# répondre
gh issue comment 95 --body "…"
gh issue close 95 --comment "Réglé par 4233568."
\end{terminal}

\section*{Lean}
\begin{terminal}
# mathlib déjà compilée
lake exe cache get
# tout compiler
lake build
# un seul fichier
lake env lean Fichier.lean
-- vérifier qu'aucun sorry ne se cache
#print axioms monTheoreme
\end{terminal}

\section*{\LaTeX{}}
\begin{terminal}
# compiler, paquets téléchargés au besoin
tectonic -X compile article.tex
\end{terminal}
